\documentclass[12pt]{article}

\usepackage[utf8]{inputenc}
\usepackage[T1]{fontenc}
\usepackage[romanian]{babel}
\usepackage{amsmath, amssymb, amsthm}
\usepackage{geometry}
\geometry{a4paper, margin=2.5cm}

\title{\Large\textbf{Teorema lui Riemann de eliminare a singularităților}}
\date{}

\theoremstyle{plain}
\newtheorem{theorem}{Teoremă}

\begin{document}

\maketitle

\begin{theorem}[Riemann]
Fie $\Omega \subset \mathbb{C}$ o mulțime deschisă și fie $a \in \Omega$.  
Presupunem că funcția


\[
f : \Omega \setminus \{a\} \to \mathbb{C}
\]


este olomorfă și este mărginită în vecinătatea punctului $a$.

Atunci $f$ admite o extensie olomorfă în întreg $\Omega$, adică există o funcție


\[
F : \Omega \to \mathbb{C}
\]


olomorfă în $\Omega$ astfel încât $F(z) = f(z)$ pentru orice $z \ne a$.

În particular, singularitatea lui $f$ în $a$ este eliminabilă.
\end{theorem}

\end{document}
